% =====================================================================
\chapter{Diszkusszió}
\label{ch:diszkusszio}
% =====================================================================

\section{Miért módszer, és nem prompt?}

A dolgozat központi állítása, hogy a vizsgagenerálás megbízhatóságát nem egy jobb prompt, hanem a
feladat explicit, ellenőrizhető lépésekre bontása biztosítja. A két vizsgálat együtt négy ponton
támasztja alá ezt az érvelést. \emph{Először}, ugyanaz a 7B modell a tervezés és az ellenőrzés hatására
a fő vizsgálat hat dokumentumán a naiv RAG $0{,}64$-es és a teljes dokumentumos prompt $0{,}83$-as
indexéről $0{,}96$-ra javult, és elérte a teljes dokumentumot kapó szervermodellt ($0{,}95$); a különbség
nem a modell tudásából, hanem a feladat szerkezetéből ered. \emph{Másodszor}, a naiv RAG minősége a
darabolás apró részleteitől függött (az angol dokumentumokon $0{,}40$-re esett), a tervezett csővezetéké
nem: a fogalmanként végzett keresés robusztusabb, mint egy feladatleírással végzett keresés.
\emph{Harmadszor}, a kérdésenkénti forráshivatkozás lehetővé tette, hogy egy olvasó -- a kiértékelésben a
bíró, a gyakorlatban az oktató -- minden kulcsot a saját forrásával szemben ellenőrizzen; éppen ez tárta
fel a naiv alapvonal hibás kulcsait, az ellenőrző karok szivárgó kérdéseit és a szervermodell forma
szerint triviális kérdéseit. \emph{Negyedszer}, ugyanaz a csővezeték a nagy szervermodellel $0{,}99$-et ért
el: a módszer nem a kis modell mankója, hanem modelltől független keret.

A vizsgálatok a módszer határait is megmutatták. Az ellenőrzés a pilotban rontott, a fő vizsgálatban mind
a hat dokumentumon javított; a fordulatot a tartalékkérdés javítása és a szivárgási szabályok hozták, vagyis
az ellenőrzés haszna az implementáció részletein múlik. Az ellenőrző ítélete a bíróéval továbbra is csak
gyengén egyezik ($\kappa\le0{,}22$): a lánc a vizsga egészét javítja, de az egyes kérdésekre adott jelzése
nem megbízható. A pilot után megfogalmazott alternatíva -- ha a más modellcsaládba tartozó ellenőrző sem
független a generátortól, a helyi telepítés helyes konfigurációja az ellenőrzés nélküli \arm{E1} -- a fő
vizsgálat alapján módosul: a javasolt helyi konfiguráció az önellenőrző \arm{E2}, a megjelölést pedig
figyelmeztetésként, nem ítéletként kell az oktatónak megmutatni. A Llama-ellenőrző (\arm{E2x}) engedékenyebb
és lassabb, a fogalomgráf (\arm{E3}) és a hibrid keresés (\arm{E2h}) nem javított, így a jelen adatokon
egyik sem éri meg a többletköltségét. Végül a teljes dokumentumot kapó szervermodell a vizsgált,
többnyire rövid dokumentumokon sem minőségben, sem lefedettségben nem maradt el: a módszer előnye vele
szemben nem a rövid dokumentumokon mért minőség, hanem az ellenőrizhetőség, az adatvédelem és a helyi
futtathatóság (P1, P3, P4), a lefedettségi előny (P2) pedig a naiv RAG-gal szemben mutatható ki.

\section{A mérés mint kutatási eredmény}

A pilot talán legfontosabb eredménye módszertani: a mérés három különböző ponton torzított volna,
ha nem validáljuk. A tartalékvizsga-hiba a generálási sikert, a csonkolt kódoló a darabolás
minőségét, a szivárgás pedig a tartalmi mutatókat és a bíró ítéletét torzította -- mindhárom
esetben \emph{felfelé}, vagyis jobb eredményt mutatva a valósnál. A három hiba közös vonása, hogy a
mérőrendszer „működött”: nem adott hibaüzenetet, és hihető számokat állított elő. Ebből az a
módszertani tanulság következik, hogy a nyelvi modellekre épülő rendszerek kiértékelésében a nyers
kimenetek kézi átolvasása és a mutatók szélsőértékeinek tudatos vizsgálata nem helyettesíthető
automatizmussal; a mi esetünkben a legjobb helyi alátámasztottsági értéket (\arm{E3}, 87\,\%) éppen
egy hibamód okozta.

A futásról futásra jelentkező zaj ($\approx0{,}07$) becslése hasonló tanulság: a kevés dokumentumon,
kevés ismétléssel mért különbségek jelentős része zaj lehet.

A fő vizsgálat a bíróval kapcsolatban két újabb vakfoltot tárt fel. A teljes dokumentumos karoknál a bíró
24\,000 karakteres ablaka a hosszú dokumentumok végéről írt helyes kérdéseket „alá nem támasztottnak”
minősíti, ami ezeket a karokat lefelé torzítja; a szervermodell visszakeresett szakaszcímre vonatkozó,
tartalmilag üres kérdéseit („hány szóból áll a cím?”) pedig a bíró hibátlannak ítélte, ami felfelé torzít.
Egy harmadik tanulság a jelentés automatizmusára vonatkozik: a jelentésgenerátor a legmagasabb indexű
kart (\arm{E4b}) „győztesnek” jelölte, holott az más dokumentumokon futott, mint a többi kar. A futásonként
eltérő dokumentumhalmazt a nézőfelület figyelmeztetésként jelzi, a rangsor azonban csak azonos
dokumentumokon futott karok között értelmes.

\section{Az érvényességet veszélyeztető tényezők}
\label{sec:validitas}

A tapasztalati szoftvertechnológiai kutatás szokásos felosztását~\cite{wohlin2012experimentation}
követve négy érvényességi dimenziót vizsgálunk.

\paragraph{Belső érvényesség.} Az egytényezős kísérleti terv a karok közötti különbséget egyetlen
tényezőhöz köti, ám a tényezők nem mindig függetlenek: az ellenőrzés például a kérdések hosszát és
formáját is megváltoztatja, ami közvetve hat a bíró érthetőségi ítéletére. A GPU-következtetés
nemdeterminisztikussága miatt azonos konfiguráció mellett is eltérő kimenetek születnek; ezt a
dokumentumszintű aggregálás és a dokumentumok számának növelése kezeli. A pilot és a fő vizsgálat
között a csővezeték változott (szivárgási szabályok, javító hívás, $r=1$), így a két vizsgálat
eredményei nem közvetlenül összevethetők; a pilot eredményei az akkori implementációra vonatkoznak. A
fő vizsgálatban az \arm{E4} és az \arm{E4b} kar más dokumentumokon futott, mint a többi, ezért ezeket csak a
közös három dokumentumon vetjük össze; a tervezett karok egyetlen maggal futottak, így a dokumentumonkénti
értékükben a futásról futásra jelentkező zaj is benne van.

\paragraph{Konstrukciós érvényesség.} A minőségi index négy egyenlő súlyú komponense közül a formai
megfelelés vizsgaszintű bináris változó, a másik három kérdésszintű ráta; ezért egyetlen formai hiba
aránytalanul nagy hatású. Ezt tudatosan vállaljuk (egy hibás formájú kérdés a gyakorlatban
használhatatlanná teszi a vizsgát), de minden komponenst külön is közlünk. A lefedettség az
embedding-térbeli legközelebbi darabra épül, ami csak közelítése annak, hogy a kérdés „miről szól”.
Az alátámasztottság és a vak megoldhatóság a bíró ítéletén alapul; a bíró a kérdés kontextusát látja,
így a dokumentum más részeiben lévő ellentmondást nem veheti észre, és -- amint a szivárgás
megmutatta -- a formailag hibás, de „igaz” kérdéseket nem bünteti; a fő vizsgálat szerint a forma
szerint triviális kérdéseket sem, a teljes dokumentumos karoknál pedig legfeljebb a dokumentum első
24\,000 karakterét látja. A bíró validálása az oktatói értékeléssel ezért elengedhetetlen; a fő vizsgálat
lezárásakor az oktatói értékelés még nem érkezett be (\tbd{kitöltendő}), így a tartalmi mutatók egyelőre
egyetlen, nem validált bíróra épülnek.

\paragraph{Külső érvényesség.} A pilot egyetlen rövid, könnyű magyar szövegen futott, amelyen a naiv
alapvonal a négy darabból hármat lát -- ez a tervezés számára a legkedvezőtlenebb helyzet. A fő vizsgálat
a 39 dokumentumból kilencen futott (karonként hat); magyar részhalmaza három rövid saját szöveg, a
Wikipédia-cikkek egyike sem szerepelt benne, így a nyelv hatása a forrás és a hossz hatásával keveredik, és
az eredmények a teljes adatkészleten módosulhatnak. A vizsgált dokumentumok többsége rövid, ezért a
hosszú dokumentumokra jellemző jelenségek (a kontextus közepének gyengébb kihasználása, a csonkolás)
csak egy-két dokumentumon jelenhettek meg. Csak
feleletválasztós kérdéseket értékelünk; az igaz--hamis és a kifejtős kérdések minősége külön vizsgálatot
igényel. Az eredmények a vizsgált modellekre (Qwen2.5-7B, Llama-3.1-8B, a szervermodell) vonatkoznak;
egy másik modellcsalád eltérően viselkedhet.

\paragraph{Következtetési érvényesség.} A dokumentumszintű nemparaméteres tesztek konzervatívak.
A fő vizsgálat karonként hat dokumentuma éppen az a legkisebb szám, amelyen $p<0{,}05$ egyáltalán
elérhető; a hat összevetés közül egy (\arm{E2}--\arm{E1}) érte el, a Holm-korrekció után egyik sem. A
fő vizsgálat következtetései ezért leíró jellegűek; a 39 dokumentum közepes hatásméretek kimutatására
elegendő lenne, kis hatások kimutatására nem. A több
összehasonlítást Holm-korrekcióval kezeljük. A bootstrap-intervallumok kis mintán alulfedhetnek, ezért a
pontbecslések mellett mindig a dokumentumszintű eloszlást is közöljük.

\section{Etikai és jogi megfontolások}

A rendszer hallgatói adatot nem kezel, a feltöltött dokumentumok a feladat végén törlődnek
(\ref{sec:zero}.~szakasz), minden generált kérdés oktatói jóváhagyást igényel, az export pedig feltünteti
az MI-közreműködést és a modellt~\cite{aiact2024}. A kiértékelő dokumentumok saját vagy nyílt licencű
(CC BY 4.0, CC BY-SA 4.0) művek; a szerverkarok és a bíró az egyetem saját GenAI-infrastruktúráján futnak.
A rendszer célja a vázlatírás gyorsítása és ellenőrizhetővé tétele, nem az oktatói döntés kiváltása.
